%% kapa: exam=1 round=30 year=2019 period=2 form=A subject=accounting subject_ko=회계학 seq=40 number=80
%% source: Q-Net 감정평가사 기출문제 (cst00309) artlSeq=5205205 file=extracted/1008/5205205/2019년 제30회 감정평가사 1차시험 문제지/2019년 제30회 감정평가사 1차시험 문제지/2019년도 제30회 감정평가사 1차 2교시 A형.hwp
%% conversion: kordoc 4.24.51 -> tools/gen1.py (auto); review=proofread-sonnet; flags=none
\begin{question}{80}
20x1년 초 영업을 개시한 상품매매기업인 (주)감평의 20x1년 1분기 월별 매출액 예산은 다음과 같다. \par
{\small\centering
\begin{tabular}{|c|c|c|c|}
\hline
 & \makecell{1월} & \makecell{2월} & \makecell{3월} \\
\hline
\makecell{매출액} & \makecell{￦2,220,000} & \makecell{￦2,520,000} & \makecell{￦2,820,000} \\
\hline
\end{tabular}\par}

(주)감평은 매출원가의 20 \%를 이익으로 가산하여 상품을 판매하고, 월말재고로 그 다음 달 매출원가의 40 \%를 보유하는 재고정책을 실시하고 있다. (주)감평의 매월 상품매입 중 50 \%는 현금매입이고, 50 \%는 외상매입이다. 외상매입대금 중 80 \%는 매입한 달의 1개월 후에, 20 \%는 매입한 달의 2개월 후에 지급된다. 상품매입과 관련하여 (주)감평의 20x1년 2월 예상되는 현금지출액은? (단, 매입에누리, 매입환출, 매입할인 등은 발생하지 않는다.) \par
\begin{choices}
  \item ￦1,076,000
  \item ￦1,100,000
  \item ￦1,345,000
  \item ￦2,176,000
  \item ￦2,445,000
\end{choices}
\end{question}
\answer{4}
